\section{Introduction}\label{sec:intro}

Two-dimensional Ti\textsubscript{3}C\textsubscript{2}T\textsubscript{x} MXene is among the most
studied electrode materials for aqueous supercapacitors~\cite{naguib2011,ghidiu2014}, and a
large body of work reports
gravimetric capacitance in H\textsubscript{2}SO\textsubscript{4} electrolyte. The
accumulated literature is an attractive target for data-driven analysis: if the reported
capacitances could be related to the structural and processing descriptors that accompany
them, the resulting model would both rank the descriptors worth optimising and supply a
data-driven prior for a subsequent physics-informed model.

That programme faces three obstacles that this study addresses directly. First, a corpus
assembled from independent publications is not an experimental design. Cell geometry,
reference electrode, mass-normalisation convention, potential window and reporting rate
differ from paper to paper, and several descriptors are reported at most once per
publication, so the unit of independent information is the paper rather than the row.
Second, such corpora are small: after discarding every entry whose capacitance is reported
only as a range or an inequality, 118 usable observations remain, drawn from 40
publications. Third, the descriptors of greatest physical interest --- interlayer spacing,
surface area, mass loading --- are exactly those most often omitted.

Synthetic data augmentation is a natural response to the second obstacle. A recent study of
flow-capacitive deionisation expanded 32 experimental observations to 200 using the
correlated-attribute (Bayesian-network) mode of \textsc{DataSynthesizer} and reported
improved model performance~\cite{fcdi_augmentation}. Whether that strategy transfers to a
corpus whose limiting factor is heterogeneity rather than row count is an open and
practically important question, and it is the question this work tests.

The contributions are as follows. (i) A leakage-controlled, paper-grouped evaluation of
gravimetric capacitance prediction on the compiled corpus. (ii) A pre-registered test of
Bayesian-network augmentation at four ratios, three network degrees and three independent
generator seeds, in which the generator is fitted inside each training fold only.
(iii) A fidelity audit establishing that the generator reproduces the training distribution
it was shown, so that the augmentation result can be interpreted rather than dismissed as a
generator failure. (iv) A variance decomposition that locates the bottleneck in
between-study heterogeneity. (v) A SHAP descriptor analysis~\cite{shap2017,treeshap2020} whose ranking is shown to be
stable under Bayesian resampling, together with an explicit assessment --- rejecting a
classical PINN~\cite{pinn2019} and specifying the physics-constrained alternative the data
do support.
